\section{舵机观测通路异步化的可行性分析}\label{sec:async-read}

\subsection{上游 TODO：基于 txPacket/rxPacket 拆分的流水线方案}\label{sec:async-read-todo}

与异步写完全没有任何上游代码痕迹相对，
HuggingFace LeRobot 主仓库的
src/lerobot/motors/motors\_bus.py
在同步读取器初始化逻辑与同步写入逻辑之间留有一段
针对“读”的 TODO 注释占位：
该注释的核心意思是：如果已经向同一组舵机和同一寄存器发出过上一轮读请求，
则下一次调用可以先接收上一轮回包，再立即发送新一轮请求；
如果请求参数发生变化，则重新配置同步读取器并退回普通同步读路径。
上游注释同时提醒，这种方法可能缩短读取等待时间，
但代价是舵机产生状态数据与策略真正使用该数据之间的延迟会增加。

注释明确指向一种“流水线（pipelining）”方案：把当前 txRxPacket 调用
中的“发请求”与“收回包”两个原本绑在一起的步骤拆开，
让每次异步读取调用先接收上一次请求的回包，
再发起这一次的新请求，下次调用再来收这一次的回包。
本质上是把单次串口往返的等待时间隐藏到两次相邻调用之间。

\subsection{与相机后台线程式异步的本质区别}\label{sec:async-read-not-thread}

需要特别澄清：上游规划的“异步读”与相机的异步读取
完全不是一个机制。
\begin{itemize}[leftmargin=2em]
  \item 相机异步读取：由独立后台线程持续调用 OpenCV 图像读取接口，
        主线程仅取最新帧。属于多线程并发模式，
        受益于每台相机独占 USB 端口，没有总线竞争问题。
  \item 舵机异步读取（TODO）：完全在主线程内执行，
        没有任何后台线程。仅仅是把单次原子的 txRxPacket 调用
        拆成两次连续调用之间的“收 + 发”重叠。属于串行流水线模式，
        因此不会与同步写入在总线上撞车。
\end{itemize}

也正因为是流水线而非并发，4.2 节中关于半双工总线竞争（约束 A）和
错误暴露延迟（约束 C）的论证对舵机异步读并不适用：
异步读不引入并发，不需要锁，错误依然在主线程当帧抛出。

\subsection{落地代价：驱动改造与观测滞后}\label{sec:async-read-cost}

尽管异步读在原理上可行，要真正落地仍有两项代价：

代价 1——需要修改 scservo\_sdk 驱动层。
当前 Feetech SDK 把 txPacket 与 rxPacket
的调用封装在 txRxPacket 里作为单个原子操作，
对外只暴露这一接口。要落地上游 TODO 的 pipelining 方案，
必须深入 scservo\_sdk 内部，把内部的 tx/rx 状态机暴露给 motors\_bus.py，
还需新增“上次发送的请求参数”缓存以及“参数不匹配则退化为同步读”的兼容路径。
这一改造涉及一个属于第三方厂商的 C/Python 混合驱动，
维护与上游升级跟踪都会有额外负担。

代价 2——观测延迟 1 帧。
TODO 注释明确写了：“This could be at the cost of increase latency between
the moment the data is produced by the motors and the moment it is used by a policy.”
异步读在第 $N$ 次调用拿到的实际上是第 $N-1$ 次请求时的电机状态。
对 ACT 这类闭环 policy 而言，
“此刻的观测”比物理时刻晚了一整帧（约 33 ms @30 FPS），
理论上会让控制误差略增。

\subsection{收益估算}\label{sec:async-read-roi}

按第三章的实测数据，单次同步读取的串口往返耗时约 1.2 ms。
异步读最多能把这 1.2 ms 从主线程的关键路径上拿掉。然而：
\begin{itemize}[leftmargin=2em]
  \item 在当前 ACT 主导的循环中（推理约 2000 ms），
        节省 1.2 ms 对总循环时间的影响小于 0.1\%；
  \item 即便后续推理被加速到 100 ms 量级（INT8 + ONNX），
        异步读省下的 1.2 ms 也仅占约 1\%，仍属边际收益；
  \item 真正可能让异步读发挥作用的场景，是在没有视觉、
        没有大模型推理的纯实时控制系统里（控制频率 500 Hz 以上），
        而这类场景与本论文的 LeRobot ACT 用例并不重合。
\end{itemize}

\subsection{小结}\label{sec:async-read-conclusion}

舵机异步读在原理上可行（pipelining 不破坏半双工约束），
但综合来看 ROI 不足以支持改造工作量：
当前推理瓶颈下省下 1.2 ms 几乎不可见，
且需要改第三方 SDK 并接受 1 帧观测滞后。
本论文将其归入“结构上可行、当前不值得做”的优化项，
作为推理加速完成后可以重新评估的潜在方向，但不在本论文的实验范围内。
